\documentclass[11pt,a4paper]{article}
\usepackage{curso}
\begin{document}
\portada{Lección}{0.4}{Repaso de álgebra lineal}

\noindent\textbf{Unidad 0.}\quad\textbf{Objetivos:} reconocer matrices SPD y su descomposición espectral;
interpretar $\kappa$; usar $e^{-At}$; aplicar Kronecker y Woodbury. \textbf{Prerrequisitos:} álgebra de
matrices; lección 0.3.

\section{¿Por qué importa?}
El solucionador lineal termodinámico (Unidad 6) relaja un sistema físico cuya dinámica es $\dot x=-(Ax-b)$ más
ruido; su tiempo de convergencia depende de $\kappa(A)$. K-FAC (Unidad 8) se apoya en productos de Kronecker y
el gradiente natural termodinámico en Woodbury. Todo se vuelve transparente con la descomposición espectral.

\section{Matrices simétricas definidas positivas}
\begin{definicion}[ (SPD)]
$A\in\R^{d\times d}$ es SPD si $A=A^\top$ y $x^\top Ax>0$ para todo $x\neq0$.
\end{definicion}
\begin{teorema}[espectral]
Si $A=A^\top$, existe $Q$ ortogonal ($Q^\top Q=I$) y $\Lambda=\diag(\lambda_i)$ real tal que $A=Q\Lambda Q^\top$.
$A$ es SPD si y solo si todos los $\lambda_i>0$.
\end{teorema}
\begin{proof}[Idea]
Para $Av=\lambda v$ con $v\neq0$ complejo: $\bar v^\top Av=\lambda\|v\|^2$ y, por simetría, $\bar v^\top Av$ es real,
luego $\lambda\in\R$. Vectores propios de valores distintos son ortogonales: $\lambda_1v_2^\top v_1=v_2^\top Av_1=
(Av_2)^\top v_1=\lambda_2v_2^\top v_1$. Si $A$ es SPD, $\lambda=v^\top Av/\|v\|^2>0$; recíprocamente
$x^\top Ax=\sum_i\lambda_i(q_i^\top x)^2>0$.
\end{proof}
Geometría: $x\mapsto Ax$ rota a la base propia ($Q^\top$), estira cada eje por $\lambda_i$ ($\Lambda$) y deshace la
rotación ($Q$). La forma cuadrática $f(x)=\frac12x^\top Ax$ es un cuenco con curvas de nivel elípticas de
semiejes $\propto1/\sqrt{\lambda_i}$ (¡ojo!: la covarianza $\Sigma$ de la lección 0.3 tiene elipses de semiejes
$\propto\sqrt{\lambda_i(\Sigma)}$, porque la densidad usa $\Sigma^{-1}$).

\section{Número de condición}
$\kappa(A)=\lambda_{\max}/\lambda_{\min}$ (para SPD; en general $\sigma_{\max}/\sigma_{\min}$). Mide (i) cuánto se
amplifican los errores relativos al resolver $Ax=b$: $\frac{\|\delta x\|}{\|x\|}\le\kappa\frac{\|\delta b\|}{\|b\|}$;
(ii) la lentitud del descenso de gradiente: con paso óptimo, el factor de contracción por iteración es
$\frac{\kappa-1}{\kappa+1}$; (iii) en un sistema físico $\dot x=-Ax$, el cociente entre la escala de tiempo más
lenta $1/\lambda_{\min}$ y la más rápida $1/\lambda_{\max}$.

\section{Exponencial de matrices}
$e^{M}=\sum_{k\ge0}M^k/k!$. La solución de $\dot x=-Ax$ es $x(t)=e^{-At}x_0$, y con $A=Q\Lambda Q^\top$:
\begin{equation}x(t)=Q\,e^{-\Lambda t}Q^\top x_0=\sum_ie^{-\lambda_it}(q_i^\top x_0)\,q_i .\end{equation}
Los modos rápidos mueren primero; la relajación total tarda $\sim1/\lambda_{\min}$ y la integración explícita
(Euler) es estable solo si $\Delta t<2/\lambda_{\max}$ (lección 2.7). El número de pasos escala como $\kappa$.
Propiedades: $e^{A+B}=e^Ae^B$ \emph{solo} si $AB=BA$; $\det e^A=e^{\Tr A}$; $(e^A)^{-1}=e^{-A}$.

\section{Productos de Kronecker}
$B\otimes C$ es la matriz por bloques $[B_{ij}C]$. Con $\mathrm{vec}$ que apila columnas:
\begin{equation}(B\otimes C)\mathrm{vec}(X)=\mathrm{vec}(CXB^\top),\quad(B\otimes C)^{-1}=B^{-1}\otimes C^{-1},\quad
(B\otimes C)(D\otimes E)=BD\otimes CE.\end{equation}
Los valores propios de $B\otimes C$ son los productos $\lambda_i(B)\mu_j(C)$. Invertir una matriz $n^2\times n^2$ cuesta
$O(n^6)$; invertir dos de $n\times n$, $O(n^3)$. K-FAC aproxima la curvatura de una capa por $A\otimes G$ (Unidad 8).

\section{Identidad de Woodbury}
\begin{resultado}
$(A+UCV)^{-1}=A^{-1}-A^{-1}U(C^{-1}+VA^{-1}U)^{-1}VA^{-1}$, con $U\in\R^{d\times k}$, $C\in\R^{k\times k}$, $V\in\R^{k\times d}$.
\end{resultado}
Caso de rango uno (Sherman--Morrison): $(\lambda I+uu^\top)^{-1}=\frac1\lambda\big(I-\frac{uu^\top}{\lambda+u^\top u}\big)$.
En aprendizaje automático la curvatura empírica $G=\lambda I+J^\top J$ con $J\in\R^{b\times d}$ y $b\ll d$ se invierte
resolviendo un sistema $b\times b$ (NGD-Woodbury, Unidad 8).

\begin{ejemplo}
$A=\begin{pmatrix}3&1\\1&2\end{pmatrix}$: $\det(A-\lambda I)=\lambda^2-5\lambda+5=0\Rightarrow\lambda=\frac{5\pm\sqrt5}2\approx3{,}618;\,1{,}382$.
Ambos positivos: SPD. $\kappa\approx2{,}618$, $\det A=5=\lambda_1\lambda_2$, $\Tr A=5=\lambda_1+\lambda_2$.
\end{ejemplo}

\begin{aplicacion}
$\kappa$ fija el tiempo del solucionador lineal termodinámico ($O(d\kappa^2\varepsilon^{-2})$ en el pensum, Unidad
6.4); $e^{-At}$ es la parte determinista del OU; Kronecker y Woodbury hacen escalable el segundo orden (Unidad 8).
\end{aplicacion}

\begin{error}
\begin{itemize}[nosep]
\item Simétrica $\not\Rightarrow$ SPD: $\begin{pmatrix}1&2\\2&1\end{pmatrix}$ tiene $\lambda=3,-1$.
\item Entradas positivas $\not\Rightarrow$ SPD (mismo ejemplo).
\item $e^{A+B}\neq e^Ae^B$ si no conmutan.
\item No formar $B\otimes C$ explícitamente: $O(n^4)$ de memoria.
\end{itemize}
\end{error}

\section{Resumen}
SPD $\Leftrightarrow\lambda_i>0$, $A=Q\Lambda Q^\top$; $\kappa=\lambda_{\max}/\lambda_{\min}$; $x(t)=e^{-At}x_0$; Kronecker y
Woodbury reducen costos de inversión.

\begin{thebibliography}{9}\small
\bibitem[30]{tla} M. Aifer et al., Thermodynamic linear algebra, \emph{npj Unconv. Comput.} 1, 13 (2024).
\bibitem[38]{kfac} J. Martens y R. Grosse, Optimizing neural networks with Kronecker-factored approximate curvature, ICML 2015.
\bibitem[39]{grosse} R. Grosse, CSC2541 lecture notes, \emph{Second-order optimization}, 2021.
\end{thebibliography}
\end{document}
